\section{An irreducible selector whose norms do not extend across every boundary}
\label{sec:r13-hensel-boundary-nonextension}

Let \(K\) be the original finite extension of \(\mathbf Q_p\),
with its original absolute value, and let \(k\) be its residue
field.  Work on the original three-dimensional torus with
coordinates \((A,C,S)\).  Let \(y\) be its unit-radius Gauss
point, so its residue field is \(k(a,c,s)\).  Put
\[
 P(V)=V^2-V-A,\qquad Z=C^2-C-A,
 \qquad F(E)=P(C+E)
       =E^2+(2C-1)E+Z.
\tag{HB1}
\]
The zero-slice divisor \(D=\{c^2-c-a=0\}\) in the smooth
residue model is smooth and irreducible.  Indeed its derivative
with respect to \(a\) is \(-1\).  Moreover \(V^2-V-a\) is
irreducible over \(k(a,s)\): a rational root, viewed as a
rational function of \(a\) over \(k(s)\), can have no finite
pole, since a pole of order \(m\) would give one of order
\(2m\) in its square minus itself.  It is therefore a
polynomial in \(a\), whose square minus itself has even
positive degree, or degree zero, and cannot equal \(a\).
The polynomial is separable in every residue characteristic.
Thus HB1 is an actual irreducible etale zero-slice presentation
at the generic point of \(D\); its derivative \(2c-1\) there
is nonzero.

Let \(\mathcal H_F\) be the strict ordinary Hensel domain
HN3, and let
\(\alpha=C+e\) be its selected root.  Since \(A\) is a
unit and the product of the two roots of \(P\) is \(-A\),
\(\alpha\) is an analytic unit on this domain.  At the
generic zero-slice point its residue is the nonzero element
\(c\in k(D)\), so this is one of the allowed finite
surface units.

\begin{proposition}[Two ordinary branches of the selected norm]
\label{prop:r13-hensel-two-ordinary-branches}
The original ordinary open domains
\[
 \mathcal B_0=\{0<|A|<1,\ 0<|C|<1\},\qquad
 \mathcal B_1=\{0<|A|<1,\ |C-1|<1\}
\tag{HB2}
\]
are contained in \(\mathcal H_F\).  At every actual point
of these domains one has
\[
             |\alpha|=|A|\quad\hbox{on }\mathcal B_0,
       \qquad |\alpha|=1\quad\hbox{on }\mathcal B_1.
\tag{HB3}
\]
No extension of the ground field or restriction to a moment
source locus occurs in these identities.
\end{proposition}

\begin{proof}
On \(\mathcal B_0\) one has \(|2C-1|=1\) and
\(|Z|<1\).  On \(\mathcal B_1\), writing \(C=1+W\),
one has \(|2C-1|=|1+2W|=1\) and
\(|Z|=|W+W^2-A|<1\).  Thus HN3 holds on both ordinary
sets, in every residue characteristic.  On the first set,
\(|C|<1\) and \(|e|<1\) imply \(|\alpha|<1\).
The equation \(\alpha^2-\alpha=A\) then gives
\(|\alpha|=|A|\).  On the second set,
\(|\alpha-1|=|C-1+e|<1\), so \(|\alpha|=1\).
\end{proof}

For \(0<r<1\), let \(q_r\) be the actual Gauss point
with radii \((r,1,1)\) in \((A,C,S)\).  For \(0<s<1\)
let \(x_0(r,s)\) be the Gauss point with radii
\((r,s,1)\) in \((A,C,S)\), and let \(x_1(r,s)\) be
the Gauss point with these same radii in \((A,C-1,S)\).
These are points of the original analytic space for all
real \(r,s\) in the displayed intervals, including
commensurable radii.  The Gauss maximum formula gives
\[
 \lim_{s\to1^-}x_0(r,s)=q_r
       =\lim_{s\to1^-}x_1(r,s),
       \qquad \lim_{r\to1^-}q_r=y.
\tag{HB4}
\]
For the first equality, the closed unit discs in the
\(C\)-variable centered at zero and at one have the
same Gauss norm; the two formulas agree on every
polynomial.  Uniform approximation on a fixed affinoid
neighborhood extends this convergence to all analytic
functions.  At \(q_r\), the two summands of \(Z\) of
degree one and two in \(C\) have independent Gauss
monomials, so \(|Z|=1\), and HN3 fails there.

\begin{theorem}[Failure of an entire central continuous extension]
\label{thm:r13-hensel-no-entire-central-continuous-extension}
Let \(U\) be any original ordinary neighborhood of
\(y\), and let \(G\) be any finite tuple of actual
analytic units on \(U\).  There is no continuous
function
\[
                     \Psi:t_G(U)\longrightarrow\mathbf R
\tag{HB5}
\]
such that
\[
             \Psi(t_G(x))=-\log|\alpha(x)|
                     \qquad(x\in U\cap\mathcal H_F).
\tag{HB6}
\]
In particular no finite actual refinement supplies a
continuous piecewise affine extension on the entire
central moment image for this selected norm.
\end{theorem}

\begin{proof}
Choose \(r<1\) sufficiently close to one that
\(q_r\in U\).  By HB4, both \(x_i(r,s)\) belong
to \(U\) for all \(s\) sufficiently close to one.
They lie in \(\mathcal B_i\), respectively.  The
actual moment map is continuous, so both moment
families converge to \(t_G(q_r)\).  Equations
HB3 and HB6 would give the two different limiting
values \(-\log r>0\) and zero for the continuous
function HB5 at this same point.  This is impossible.
The argument applies to every finite tuple \(G\),
not only to the coordinates used to exhibit HB4.
\end{proof}

All coefficients of \(P\) have norm at most one
at \(y\), and
\[
                         |\operatorname{Res}(P,V)(y)|
                              =|-A(y)|=1.
\tag{HB7}
\]
Thus the all-conjugate resultant hypotheses of HN14
hold.  The selected logarithm nevertheless lacks the
entire extension in HB5.  There is no contradiction:
its value tends to zero along every approach to
\(y\) in the selected domain, whereas HB4 gives
two different limits at each nearby excluded point
\(q_r\).  Protected continuity on the realized
selected domain, and continuity after adjoining the
single original center, are the precise surviving
statements.  This example refutes the stronger
extension shortcut; it does not obstruct an intrinsic
differential construction using actual local root
domains or a separately justified singular pairing.
